\section{Background and related work}

\subsection{The space link and its security services}

Telecommand and telemetry in most civil and commercial missions follow the CCSDS
protocol stack: application data in Space Packets \cite{ccsds133spp}, carried in
transfer frames, optionally protected by the Space Data Link Security protocol
(SDLS) \cite{ccsds355sdls}. SDLS provides authentication, optional confidentiality,
and --- critically for this work --- an anti-replay sequence number bound into the
authenticated data. Our testbed models SDLS at the level of abstraction the
detections care about: a keyed authentication tag over the frame payload and
sequence number, and a sliding acceptance window on the vehicle. It does not
model bit layouts, which do not change which events reach the analytics tier.

The two properties that matter for detection engineering are worth stating
plainly. First, a frame that fails its tag check on the vehicle is an
unambiguous signal: nominal operations never produce one, because the ground
system signs every frame with the key the vehicle holds. Second, authentication
alone does not establish freshness --- a recorded frame remains cryptographically
valid forever --- which is why the anti-replay window, not the tag, is the
control that defeats replay.

\subsection{Threat knowledge for space systems}

SPARTA \cite{sparta} adapts the ATT\&CK \cite{attack} structure to space systems,
covering reconnaissance through impact for both the space and ground segments. We
map every detection rule to a SPARTA technique and, where an analogue exists, to
an ATT\&CK technique; this makes the coverage claim checkable rather than
rhetorical. NIST IR 8270 \cite{nistir8270} and SPD-5 \cite{spd5} supply the
control expectations --- authentication of commands, protection of telemetry
integrity, monitoring --- that the architecture implements.

\subsection{Prior empirical work}

Pavur and Martinovic \cite{pavur2020launchpad} argue that satellite security
research is held back by the absence of accessible testbeds, and that the field
substitutes speculation for measurement as a result. Willbold et al.
\cite{willbold2023odyssey} answer part of that challenge on the space segment,
analysing real satellite firmware and finding that several vehicles accept
unauthenticated commands outright. Manulis et al. \cite{manulis2021newspace}
survey the New Space threat landscape and identify the ground segment as the
dominant practical attack surface, a judgement the KA-SAT incident
\cite{viasat2022} supports.

Our work is complementary and deliberately narrower. We do not analyse flight
software. We assume the space segment implements its security services correctly
and ask what the \emph{ground} segment's cloud telemetry sees when an adversary
attacks the chain around it. To our knowledge this combination --- an end-to-end
ground-segment testbed, five labelled attack scenarios, and reported detection
latency, false-positive rate, control attribution and cost --- has not previously
been published in a form a reader can rerun.

\subsection{Why detection latency is the right metric here}

In enterprise security, dwell time is measured in days and the marginal value of
shaving seconds is low. A ground segment inverts this. The adversary's window of
\emph{effect} is the pass: a command that cannot be transmitted cannot harm the
vehicle. A pass in our configuration lasts about \SI{7}{\minute}. A detection
that fires in seconds can contain an intrusion before the next pass; one that
fires in an hour cannot. Latency is therefore not a quality-of-service metric in
this domain --- it decides whether the response is capable of mattering.
